@tex
%% Cygnus revisions to texinfo.tex, TeX macros to handle texinfo files
%   Copyright (C) 1991 Free Software Foundation, Inc.
%% Maintained at Cygnus Support as: 
%%$Id: texiplus.tex,v 1.14 1991/12/06 23:38:51 pesch Exp $
  
%This texi+.tex file is free software; you can redistribute it and/or
%modify it under the terms of the GNU General Public License as
%published by the Free Software Foundation; either version 1, or (at

{\def\$#1${{ #1}} % Kluge: collect RCS revision info without $...$
\let\fsfvn=\texinfoversion
\xdef\texinfoversion{\fsfvn\ Cygnus \$Revision: 1.14 $}   % Version filled in by RCS.
}
\write16{Loading Cygnus texinfo revisions [\texinfoversion]:}
  
% Print the version number if in a .fmt file.
\everyjob{\write16{[Cygnus Texinfo \texinfoversion]}}

\globaldefs=1 % Escape the bounds of @tex/@end tex surrounding us

%%%%%%%%%%%%%%%%%%%%%%%%%%%%%%%CROPMARKS%%%%%%%%%%%%%%%%%%%%%%%%%%%%%%
% These differ only slightly from FSF defaults; all crop dimens are
% defined here (whether different or not) for ease in revising.
\cornerlong=1pc\cornerthick=.3pt	% These set size of cropmarks
\outerhsize=7in
\outervsize=9in
%%%%%%%%%%%%%%%%%%%%%%%%%%%%%%%%%%%%%%%%%%%%%%%%%%%%%%%%%%%%%%%%%%%%%%
%%%%%%%%%%%%%%%%%%%%%%%%%CROPPAGEOUT%%%%%%%%%%%%%%%%%%%%%%%%%%%%%%%%%%
% In the long term we DO NOT need our own copy of \croppageout.  At
% the moment, though, there's a debate about whether texinfo.tex's
% version should stop setting \hoffset=0.  We want that NOW and we don't
% use \bindingoffset, so for the nonce put a slightly modified
% \croppageout here.  The only intended change from the texinfo.tex
% version is to avoid setting \hoffset.
\def\croppageout#1{\shipout
		 \vbox to \outervsize{\hsize=\outerhsize
                 \vbox{\line{\ewtop\hfill\ewtop}}
                 \nointerlineskip
                 \line{\vbox{\moveleft\cornerthick\nstop}
                       \hfill
                       \vbox{\moveright\cornerthick\nstop}}
                 \vskip \topandbottommargin
                 \centerline{\ifodd\pageno\hskip\bindingoffset\fi
			\vbox{
			{\let\hsize=\pagewidth \makeheadline}
			\pagebody{#1}
			{\let\hsize=\pagewidth \makefootline}}
			\ifodd\pageno\else\hskip\bindingoffset\fi}
		 \vskip \topandbottommargin plus1fill minus1fill
                 \boxmaxdepth\cornerthick
                 \line{\vbox{\moveleft\cornerthick\nsbot}
                       \hfill
                       \vbox{\moveright\cornerthick\nsbot}}
                 \nointerlineskip
                 \vbox{\line{\ewbot\hfill\ewbot}}
	}
  \advancepageno 
  \ifnum\outputpenalty>-20000 \else\dosupereject\fi}
%
%%%%%%%%%%%%%%%%%%%%%%%%%%%%%%%%%%%%%%%%%%%%%%%%%%END CROPPAGEOUT
		
\newif\ifdraft\drafttrue % Extra markings; turn off with @finalout

% There is only a small change to standard \title from texinfo 
% ...to include DRAFT marking on title page unless @finalout.
% However, since \title is local to \titlepage, we have to redefine
% *that* whole damned thing.  
%%%%%%%%%%This is sensitive to conflict w/FSF changes!****************
\def\titlepage{\begingroup \parindent=0pt \textfonts
   \let\subtitlerm=\tenrm
% I deinstalled the following change because \cmr12 is undefined.
% This change was not in the ChangeLog anyway.  --rms.
%   \let\subtitlerm=\cmr12
   \def\subtitlefont{\subtitlerm \normalbaselineskip = 13pt \normalbaselines}%
   %
   \def\authorfont{\authorrm \normalbaselineskip = 16pt \normalbaselines}%
   %
   % Leave some space at the very top of the page.
   \vglue\titlepagetopglue
   %
   % Now you can print the title using @title.
   \def\title{\parsearg\titlezzz}%
%change from FSF only in following line:
   \def\titlezzz##1{\leftline{\titlefont{##1\ifdraft\hfill DRAFT\fi }}
%end real change from FSF version
		    % print a rule at the page bottom also.
		    \finishedtitlepagefalse
		    \vskip4pt \hrule height 4pt \vskip4pt}%
   % No rule at page bottom unless we print one at the top with @title.
   \finishedtitlepagetrue
   %
   % Now you can put text using @subtitle.
   \def\subtitle{\parsearg\subtitlezzz}%
   \def\subtitlezzz##1{{\subtitlefont \rightline{##1}}}%
   %
   % @author should come last, but may come many times.
   \def\author{\parsearg\authorzzz}%
   \def\authorzzz##1{\ifseenauthor\else\vskip 0pt plus 1filll\seenauthortrue\fi
      {\authorfont \leftline{##1}}}%
   %  
   % Most title ``pages'' are actually two pages long, with space
   % at the top of the second.  We don't want the ragged left on the second.
   \let\oldpage = \page
   \def\page{%
      \iffinishedtitlepage\else
	 \finishtitlepage
      \fi
      \oldpage
      \let\page = \oldpage
      \hbox{}}%
%   \def\page{\oldpage \hbox{}}
}%%%%%%%%%%end sensitive to conflict w/FSF changes!********************

% Cygnus uses different default page headings and footers in
% singlepage and doublepage layouts:
\def\leadline{\ \leaders\hrule height 2.5pt depth -2pt\hfil\ }
% SINGLEPAGE:
\def\HEADINGSsingle{
\global\pageno=1
\evenfootline={\line{\let\,=\thinspace\ftfnt c\,y\,g\,n\,u\,s\quad s\,u\,p\,p\,o\,r\,t\leadline\manvers}}
\oddfootline={\line {\let\,=\thinspace\ftfnt c\,y\,g\,n\,u\,s\quad s\,u\,p\,p\,o\,r\,t\leadline\manvers}}
\global\evenheadline={\line{\ftfnt\thischapter\leadline
		\ifdraft DRAFT\leadline\fi
		\bf\folio}}
\global\oddheadline={\line{\ftfnt\thischapter\leadline
		\ifdraft DRAFT\leadline\fi
		\bf\folio}}
}
%
% DOUBLEPAGE:
% use hrule leaders to delimit headings, footings from
%               body; timestamp footer; pagenos outside *bottom* to
%               permit moving chaptername to outside *top* (for easier
%               skimming).  1990 dec 31, pesch@cygnus.com
\def\manvers{}% Empty default version in case manual doesn't supply
\def\HEADINGSdouble{
\pageno=1
\evenfootline={\vbox{\line{{\bf\folio}\ftfnt\leadline\ifdraft DRAFT\quad\fi\ftfnt\manvers\qquad\today}}}
\oddfootline={\vbox{\line{{\let\,=\thinspace\ftfnt c\,y\,g\,n\,u\,s\quad s\,u\,p\,p\,o\,r\,t}\leadline{\bf\folio}}}}
\evenheadline={\vbox{\line{\ftfnt\thistitle\leadline\ifdraft DRAFT\fi}}}
\oddheadline={\vbox{\line{\rm\leadline\thischapter}}}
}

% CHANGE to @finalout ---also use it to remove DRAFT markings from
% title, footers
%% For a final copy, take out the rectangles
%% that mark overfull boxes (in case you have decided
%% that the text looks ok even though it passes the margin).
\def\finalout{\overfullrule=0pt
%% Also take out ``DRAFT'' markings
	\global\draftfalse
			}

%Font overrides; we use PostScript when possible.
\font\textrm=pncr scaled \magstephalf
\font\texttt=pcrr
\font\textbf=pncb scaled \magstephalf 
\font\textit=pncri scaled \magstephalf
\font\textsl=pcrro
\font\textsf=slcrst at 9.5pt
% FIXME: figure out, test dvips smallcaps kluge
%\font\textsc=Times-SmallCaps scaled \magstephalf
\font\ftfnt=phvr at 8pt % Cygnus Support footer-font

\font\defbf=pncb scaled \magstep1 %was 1314

\font\ninett=pcrr at 9pt
\let\indtt=\ninett
\font\indrm=pncr at 9pt
\font\indit=pncri at 9pt
\font\indsl=pcrro at 9pt
\let\indsf=\indrm
\let\indbf=\indrm
\let\indsc=\indrm

\font\chaprm=pncb at 17pt
\font\chapit=pncbi at 17pt
\font\chapsl=pcrbo at 16pt
\font\chaptt=pcrb at 16pt
\font\chapsf=slcrst at 16pt
\let\chapbf=\chaprm

\font\secrm=pncb at 14pt
\font\secit=pncbi at 14pt
\font\secsl=pcrbo at 13pt
\font\sectt=pcrb at 13pt
\font\secsf=slcrst at 13pt
\font\secbf=pncb at 14pt

\font\ssecrm=pncr at 13pt
\font\ssecit=pncri at 13pt
\font\ssecsl=pcrro at 12pt
\font\ssectt=pcrr at 12pt
\font\ssecsf=slcrst at 12pt
\font\ssecbf=pncb at 13pt

\font\titlerm=pncb at 21pt
\font\authorrm=pncr scaled \magstep2

\font\truesecrm=pncr at 12pt

% Fonts for short table of contents
\font\shortcontrm=pncr at 12pt 
\font\shortcontbf=pncb at 12pt 
\font\shortcontsl=pncri at 12pt

% \smartitalic{ARG} outputs arg in italics, followed by an italic correction
% unless the following character is such as not to need one.
\def\smartitalicx{\ifx\next,\else\ifx\next-\else\ifx\next.\else\/\fi\fi\fi}
\def\smartitalic#1{{\it #1}\futurelet\next\smartitalicx}
% \smartitalic is used for @emph; FSF keeps trying to make it match
% @var.
%
% \smartslant{ARG} outputs arg in a slanted font, followed by an
% italic correction unless the following character is such as not to need one.
\def\smartslantx{\ifx\next,\else\ifx\next-\else\ifx\next.\else\/\fi\fi\fi}
\def\smartslant#1{{\sl #1}\futurelet\next\smartslantx}

\let\i=\smartitalic
\let\var=\smartslant
\let\dfn=\smartitalic
\let\emph=\smartitalic
\let\cite=\smartitalic

%%%%%%%%%%%%%%%%%%%%%%%%%%%%%%%%%%%%%%%%%%%%%%%%%%%%%%%%%%%%%%%%%%%%%%
% @altref, if called explicitly (inside @iftex), switches the @ref
% printed output so it lacks the square brackets used by default.
% Unfortunately this requires reproducing the entire damn xrefX
% definition. WARNING: may need to track texinfo.tex changes to xrefX...
\newif\ifbra\bratrue
\def\altref{\ifbra\brafalse\else\bratrue\fi}% Toggle.
%%%%%%%%%%%%%%%%%%%%%%%%%%%%%%Begin duplicate xrefX:
\def\xrefX[#1,#2,#3,#4,#5,#6]{\begingroup%
\def\printedmanual{\ignorespaces #5}%
\def\printednodename{\ignorespaces #3}%
%
\setbox1=\hbox{\printedmanual}%
\setbox0=\hbox{\printednodename}%
\ifdim \wd0=0pt%
\def\printednodename{\ignorespaces #1}%
%%% Uncommment the following line to make the actual chapter or section title
%%% appear inside the square brackets.
%\def\printednodename{#1-title}%
\fi%
%
%
% If we use \unhbox0 and \unhbox1 to print the node names, TeX does
% not insert empty discretionaries after hyphens, which means that it
% will not find a line break at a hyphen in a node names.  Since some
% manuals are best written with fairly long node names, containing
% hyphens, this is a loss.  Therefore, we simply give the text of
% the node name again, so it is as if TeX is seeing it for the first
% time.
\ifdim \wd1>0pt
section ``\printednodename'' in \cite{\printedmanual}%
\else%
\turnoffactive%
%%%%%%%%%%%%%%%Here's the change for @altref:
\ifbra
\refx{#1-snt}{} ``\printednodename'', page\tie\refx{#1-pg}{}%
\else
\refx{#1-snt}{} \printednodename, page\tie\refx{#1-pg}{}%
\fi
%%%%%%%%%%%%%%%end change for @altref
\fi
\endgroup}
%%%%%%%%%%%%%%%%%%%%%%%%%%%%%%%:End duplicate xrefX
%%%%%%%%%%%%%%%%%%%%%%%%%%%%%%%%%%%%%%%%%%%%%%%%%%%%%%%%%%%%
% \widen: let examples (or whatever) fall into margins.  Symmetric.  Cancel
% ``@widen{N}'' with ``@widen{-N}''.
\newdimen\extra\extra=0pt % for ``widen''
\global\def\widen#1{\advance\extra by #1%
\advance\leftskip by -#1\advance\rightskip by -#1}
%
%%%%%%%%%%%%%%%%%%%%%%%%%%%%%%%%%%%%%%%%%%%%%%%%%%%%%%%%%%%%
%%%%%%%%%%%%%%%%%%%%%%%%%%%%%%%%%%%%%%%%%%%%%%%%%%MODIFIED \cartouche
% Only intended change from texinfo.tex is \lispnarrowing turned off
% within cartouche.
\long\def\cartouche{%
\begingroup
	\lispnarrowing=0in % cartouche serves as alternate highlighting
	\lskip=\leftskip \rskip=\rightskip
	\leftskip=0pt\rightskip=0pt %we want these *outside*.
	\cartinner=\hsize \advance\cartinner by-\lskip 
		 	  \advance\cartinner by-\rskip
	\cartouter=\hsize
	\advance\cartouter by 18pt % allow for 3pt kerns on either
%				     side, and for 6pt waste from
%				     each corner char
	\normbskip=\baselineskip \normpskip=\parskip \normlskip=\lineskip
	\vbox\bgroup
		\baselineskip=0pt\parskip=0pt\lineskip=0pt
		\carttop
		\hbox\bgroup
			\hskip\lskip
			\vrule\kern3pt
			\vbox\bgroup
				\hsize=\cartinner
				\kern3pt
				\begingroup
					\baselineskip=\normbskip
					\lineskip=\normlskip
					\parskip=\normpskip
					\vskip -\parskip
\def\Ecartouche{%
				\endgroup
				\kern3pt
			\egroup
			\kern3pt\vrule
			\hskip\rskip
		\egroup
		\cartbot
	\egroup
\endgroup
}}	
%%%%%%%%%%%%%%%%%%%%%%%%%%%%%%%%%%%%%%%%%%%%%%%%END MODIFIED \cartouche
\textfonts

\globaldefs=0 % this is NOT redundant; the \endgroup done by@end tex
	      % would restore the value, but before doing that would
	      % screw up if globaldefs=1 were allowed to remain here.
@end tex
@rm
